\section*{Capítulo \ref{ch:g9:metals-acids-corrosion} --- Os metais: reações com os ácidos e corrosão}

\begin{solution}{exo:g9:metals-acids-corrosion:1}
O hidrogênio: uma tala acesa na boca do tubo dá um estalo agudo.
\end{solution}

\begin{solution}{exo:g9:metals-acids-corrosion:2}
O ferro, o zinco e o alumínio reagem; o cobre e o ouro não.
\end{solution}

\begin{solution}{exo:g9:metals-acids-corrosion:3}
O oxigênio e a água, ao mesmo tempo.
\end{solution}

\begin{solution}{exo:g9:metals-acids-corrosion:4}
Um \omterm{def:g9:ions:ion}{íon} presente durante uma reação, que não participa dela. Com o ácido
clorídrico: o \omterm{def:g9:ions:ion}{íon} cloreto \ce{Cl-}.
\end{solution}

\begin{solution}{exo:g9:metals-acids-corrosion:5}
\ce{Mg(s) + 2H+(aq) -> Mg^{2+}(aq) + H2(g)}; carga $+2$ de cada lado.
\end{solution}

\begin{solution}{exo:g9:metals-acids-corrosion:6}
Acrescentar hidróxido de sódio: um \omterm{def:g9:ions:precipitate}{precipitado} verde de \ce{Fe(OH)2}
mostra os \omterm{def:g9:ions:ion}{íons} ferro(II).
\end{solution}

\begin{solution}{exo:g9:metals-acids-corrosion:7}
A fervura retira o \omterm{def:g4:air-a-mixture-of-gases:air}{ar} dissolvido na água, e o óleo impede que o \omterm{def:g4:air-a-mixture-of-gases:air}{ar} se
dissolva de novo: o prego tem água, mas não tem oxigênio.
\end{solution}

\begin{solution}{exo:g9:metals-acids-corrosion:8}
O sal dissolvido na água que respinga no carro acelera a \omterm{def:g9:metals-acids-corrosion:corrosion}{ferrugem}.
\end{solution}

\begin{solution}{exo:g9:metals-acids-corrosion:9}
A camada de óxido formada gruda no metal e o isola do \omterm{def:g4:air-a-mixture-of-gases:air}{ar}: o ataque para
na superfície.
\end{solution}

\begin{solution}{exo:g9:metals-acids-corrosion:10}
Pintá-lo, lubrificá-lo com óleo ou graxa, galvanizá-lo (ou fabricá-lo de aço inoxidável).
\end{solution}

\begin{solution}{exo:g9:metals-acids-corrosion:11}
\ce{2Al(s) + 6H+(aq) -> 2Al^{3+}(aq) + 3H2(g)}: 2 Al e 6 H de cada lado,
carga $+6$ de cada lado. Para 200 \omterm{def:g7:atoms-and-molecules:atom}{átomos} de Al: $200 \times \frac{3}{2} = 300$
\omterm{def:g7:atoms-and-molecules:molecule}{moléculas} de hidrogênio.
\end{solution}

\begin{solution}{exo:g9:metals-acids-corrosion:12}
O zinco em volta do arranhão é corroído no lugar do aço: enquanto
houver zinco perto, o aço é poupado.
\end{solution}

\begin{solution}{pb:g9:metals-acids-corrosion:1}
\textbf{1.} \ce{Zn(s) + 2H+(aq) -> Zn^{2+}(aq) + H2(g)}.

\textbf{2.} O hidrogênio: um estalo agudo com uma tala acesa.

\textbf{3.} O hidróxido de sódio dá um \omterm{def:g9:ions:precipitate}{precipitado} branco de
\ce{Zn(OH)2}.

\textbf{4.} Os \omterm{def:g9:ions:ion}{íons} cloreto \ce{Cl-}.

\textbf{5.} $125.6 - 113.5 = \qty{12.1}{g}$.

\textbf{6.} O hidrogênio é liberado enquanto o zinco reage com o ácido;
quando não sobra mais zinco, o borbulhamento para.

\textbf{7.} O ferro do aço também reagiria, dando hidrogênio e \omterm{def:g9:ions:ion}{íons}
\ce{Fe^{2+}}; o hidróxido de sódio daria então um \omterm{def:g9:ions:precipitate}{precipitado} verde.

\textbf{8.} O zinco é atacado no lugar do aço e mantém a água e o \omterm{def:g4:air-a-mixture-of-gases:air}{ar}
longe dele.

\textbf{9.} $2 \times 10.0 \times 10.0 = \qty{200}{cm^2}$.

\textbf{10.} $12.1 \div 7.13 \approx \qty{1.70}{cm^3}$.

\textbf{11.} $1.70 \div 200 = \qty{0.0085}{cm}$.

\textbf{12.} $0.0085 \times 10\,000 = \qty{85}{\micro m}$: a camada de
zinco tinha cerca de \qty{85}{\micro m} de espessura.
\end{solution}
